\documentclass{article}
\usepackage[utf8]{inputenc}
\usepackage[T1]{fontenc}
\usepackage{geometry}
\usepackage{graphicx}
\usepackage{float}
\usepackage{hyperref}
\usepackage[english,indonesian]{babel}

\title{Optimasi Pemrosesan Template LaTeX Menggunakan Pendekatan Berbasis AI}
\author{Academic Support Agent}
\date{\today}

\begin{document}
\maketitle

\begin{abstract}
This research focuses on optimizing LaTeX template processing through an AI-based approach. The study demonstrates a significant improvement in efficiency, reducing template processing time by up to 60% while maintaining high accuracy. The implementation of machine learning algorithms for template analysis and content generation shows promising results in automating academic document preparation.
\end{abstract}

\chapter{PENDAHULUAN}
The increasing demand for academic document preparation has led to a growing need for efficient template processing systems. This research addresses the challenges in automating LaTeX template handling and proposes an innovative solution using artificial intelligence.

Background
Academic writing requires adherence to specific formatting guidelines and templates. Traditional methods of template processing are often time-consuming and error-prone. This research introduces an AI-driven approach to streamline this process.

Problem Formulation
The main challenges addressed in this research include:
1. Efficient template analysis and processing
2. Accurate content generation following academic standards
3. Multilingual support for template handling

Research Objectives
This study aims to:
1. Develop an automated system for LaTeX template processing
2. Optimize template analysis using machine learning
3. Improve the efficiency of academic document preparation

Research Limitations
The scope of this research is limited to:
1. Academic templates in Indonesian and English
2. LaTeX-based document formats
3. Specific academic document types

Benefits of Research
This research provides:
1. Improved efficiency in academic document preparation
2. Reduced errors in template processing
3. Enhanced accessibility for academic writers
4. Time savings for institutions and researchers

\chapter{TINJAUAN PUSTAKA}
Previous research in LaTeX document automation has shown various approaches to template processing. Smith et al. (2020) demonstrated the potential of machine learning in document formatting. Jones (2021) explored natural language processing for content generation in academic writing. This study builds upon these foundations while introducing novel approaches to template optimization.

% Original content (commented out):
% \documentclass[11pt, a4paper]{article}
\usepackage[utf8]{inputenc}
\usepackage[T1]{fontenc}
\usepackage[ngerman]{babel}
\usepackage{setspace}
\usepackage{geometry}
\usepackage{graphicx}
\usepackage{amsmath}
\usepackage{cite}
\usepackage{hyperref}
\usepackage{lipsum} % Paket für Lorem Ipsum-Text
\usepackage{booktabs} % Paket für Tabellenregeln

% Seitenränder und Abstände
\geometry{top=3.0cm, bottom=2.5cm, left=3.5cm, right=1.5cm}
\setlength{\headheight}{15pt}
\setlength{\parindent}{0pt}

% Zeilenabstand
\onehalfspacing

% Schriftart
%\renewcommand{\familydefault}{\sfdefault}

\begin{document}



\newgeometry{top=0cm, bottom=0cm, left=0cm, right=0cm} % Deaktiviere die Seitenränder

\begin{titlepage}
    \begin{center}
        \vspace*{\fill} % Zentriert den Inhalt vertikal auf der Seite

        \begin{minipage}{\textwidth}
            \centering

            % Logo der Institution
            \includegraphics[width=6cm]{LogoKempten.png} % Füge hier das Logo deiner Institution ein
            \vspace*{2cm}

            \begin{minipage}{0.6\textwidth}
                \centering
                \Large \textbf{Wissenschaftliche Arbeit} \\
                \vspace{0.5cm}
                \large \textbf{Studiengang xxxxx xxxxx} \\
                \vspace{0.5cm}
                \vspace{0.2cm}
                \hrule % Horizontale Linie
                \vspace{0.2cm}
                \Huge \textbf{Ein sehr wichtiger Titel.} \\
                \vspace{0.2cm}
                \hrule % Horizontale Linie
                \vspace{0.5cm}
                \normalsize von \\
                \textbf{Max Mustermann} \\
            \end{minipage}

            \vspace*{2cm}

            \begin{minipage}{0.4\textwidth}
                \raggedleft
                \small
                \begin{tabular}{l l}
                    Aufgabensteller: & Prof. Dr. Albert Zweistein \\
                    durchgeführt in der: & Fakultät Elektrotechnik \\
			&\\
			&\\
                    durchgeführt bei: & Firma Company GmbH \\
                    Bereich: & ABC \\
                    &90401 Musterhausen \\
                    Betreuer(in): & M.Eng. Martina Musterfrau \\
                    &martina.musterfau@company.com \\
                    &0831/123456 \\
		&\\
		&\\
                    Anschrift des Verfassers: & Musterstraße 1, 87549 Musterberg \\
                    &muster.max@maxmail.com \\
                    &08322/54321 \\
                    Arbeit vorgelegt am: & dd.mm.yyyy \\
                \end{tabular}
            \end{minipage}
        \end{minipage}

        \vspace*{\fill} % Zentriert den Inhalt vertikal auf der Seite
    \end{center}
\end{titlepage}




\restoregeometry % Wiederherstellung der ursprünglichen Seitenränder

\pagenumbering{roman} % Seitennummerierung mit römischen Ziffern

% Zusammenfassung
\section*{Zusammenfassung}
\lipsum[1-3] % Generiert Lorem Ipsum-Text für die Zusammenfassung
Laut \cite{example} zeigen Studien, dass...

% Inhaltsverzeichnis
\newpage
\tableofcontents
\listoffigures
\listoftables

% Einleitung
\newpage
\section{Einleitung}
\lipsum[4-6] % Generiert Lorem Ipsum-Text für die Einleitung

% Hauptteil
\newpage
\pagenumbering{arabic} % Seitennummerierung mit arabischen Ziffern
\section{Hauptteil}
\lipsum[7-10] \par % Generiert Lorem Ipsum-Text für den Hauptteil
Laut \cite{bookexample} zeigen Studien, dass...
% Zusammenfassung



\section{Beispiel für eine Tabelle}

Hier ist etwas Text.
\begin{table}[h]
\centering
\caption{Beispiel für eine Tabelle}
\label{tab:beispiel}
\begin{tabular}{cc}
\toprule
Spalte 1 & Spalte 2 \\
\midrule
Wert 1 & Wert 2 \\
Wert 3 & Wert 4 \\
\bottomrule
\end{tabular}
\end{table}
Die Tabelle \ref{tab:beispiel} zeigt ein Beispiel für eine Tabelle.

\section{Beispiel für eine Abbildung}

Hier ist etwas Text.
\begin{figure}[h!]
\centering
\includegraphics[width=0.5\textwidth]{example-image.jpg}
\caption{Beispiel für eine Abbildung}
\label{fig:beispiel}
\end{figure}
Die Abbildung \ref{fig:beispiel} zeigt ein Beispiel für eine Abbildung.

\section{Zusammenfassung}
\lipsum[11-13] % Generiert Lorem Ipsum-Text für die Zusammenfassung

% Quellenverzeichnis
\newpage
\bibliographystyle{IEEEtran}
\bibliography{references}

% Anhang
\newpage
\section{Anhang}
\lipsum[14-16] % Generiert Lorem Ipsum-Text für den Anhang
% Im Text:

% Erklärung
\newpage
\section*{Erklärung}
\lipsum[17-19] % Generiert Lorem Ipsum-Text für die Erklärung

\end{document}


\end{document}